%-----------------------------------------------------------------------------------------------------------------------------------------------%
%	The MIT License (MIT)
%
%	Copyright (c) 2021 Jitin Nair
%
%	Permission is hereby granted, free of charge, to any person obtaining a copy
%	of this software and associated documentation files (the "Software"), to deal
%	in the Software without restriction, including without limitation the rights
%	to use, copy, modify, merge, publish, distribute, sublicense, and/or sell
%	copies of the Software, and to permit persons to whom the Software is
%	furnished to do so, subject to the following conditions:
%	
%	THE SOFTWARE IS PROVIDED "AS IS", WITHOUT WARRANTY OF ANY KIND, EXPRESS OR
%	IMPLIED, INCLUDING BUT NOT LIMITED TO THE WARRANTIES OF MERCHANTABILITY,
%	FITNESS FOR A PARTICULAR PURPOSE AND NONINFRINGEMENT. IN NO EVENT SHALL THE
%	AUTHORS OR COPYRIGHT HOLDERS BE LIABLE FOR ANY CLAIM, DAMAGES OR OTHER
%	LIABILITY, WHETHER IN AN ACTION OF CONTRACT, TORT OR OTHERWISE, ARISING FROM,
%	OUT OF OR IN CONNECTION WITH THE SOFTWARE OR THE USE OR OTHER DEALINGS IN
%	THE SOFTWARE.
%	
%
%-----------------------------------------------------------------------------------------------------------------------------------------------%

%----------------------------------------------------------------------------------------
%	DOCUMENT DEFINITION
%----------------------------------------------------------------------------------------

% article class because we want to fully customize the page and not use a cv template
\documentclass[a4paper,12pt]{article}

%----------------------------------------------------------------------------------------
%	FONT
%----------------------------------------------------------------------------------------

% % fontspec allows you to use TTF/OTF fonts directly
% \usepackage{fontspec}
% \defaultfontfeatures{Ligatures=TeX}

% % modified for ShareLaTeX use
% \setmainfont[
% SmallCapsFont = Fontin-SmallCaps.otf,
% BoldFont = Fontin-Bold.otf,
% ItalicFont = Fontin-Italic.otf
% ]
% {Fontin.otf}

%----------------------------------------------------------------------------------------
%	PACKAGES
%----------------------------------------------------------------------------------------
\usepackage{url}
\usepackage{parskip} 	

%other packages for formatting
\RequirePackage{color}
\RequirePackage{graphicx}
\usepackage[usenames,dvipsnames]{xcolor}
\usepackage[scale=0.9]{geometry}

%tabularx environment
\usepackage{tabularx}

%for lists within experience section
\usepackage{enumitem}

% centered version of 'X' col. type
\newcolumntype{C}{>{\centering\arraybackslash}X} 

%to prevent spillover of tabular into next pages
\usepackage{supertabular}
\usepackage{tabularx}
\newlength{\fullcollw}
\setlength{\fullcollw}{0.47\textwidth}

%custom \section
\usepackage{titlesec}				
\usepackage{multicol}
\usepackage{multirow}

%CV Sections inspired by: 
%http://stefano.italians.nl/archives/26
\titleformat{\section}{\Large\scshape\raggedright}{}{0em}{}[\titlerule]
\titlespacing{\section}{0pt}{10pt}{10pt}

%for publications
\usepackage[style=authoryear,sorting=ynt, minbibnames=2, maxbibnames=2]{biblatex}

%Setup hyperref package, and colours for links
\usepackage[unicode, draft=false]{hyperref}
\definecolor{linkcolour}{rgb}{0,0.2,0.6}
\hypersetup{colorlinks,breaklinks,urlcolor=linkcolour,linkcolor=linkcolour}
\addbibresource{citations.bib}
\setlength\bibitemsep{1em}

%for social icons
\usepackage{fontawesome5}

%debug page outer frames
%\usepackage{showframe}


% job listing environments
\newenvironment{jobshort}[2]
    {
    \begin{tabularx}{\linewidth}{@{}l X r@{}}
    \textbf{#1} & \hfill &  #2 \\[3.75pt]
    \end{tabularx}
    }
    {
    }

\newenvironment{joblong}[2]
    {
    \begin{tabularx}{\linewidth}{@{}l X r@{}}
    \textbf{#1} & \hfill &  #2 \\[3.75pt]
    \end{tabularx}
    \begin{minipage}[t]{\linewidth}
    \begin{itemize}[nosep,after=\strut, leftmargin=1em, itemsep=3pt,label=--]
    }
    {
    \end{itemize}
    \end{minipage}    
    }



%----------------------------------------------------------------------------------------
%	BEGIN DOCUMENT
%----------------------------------------------------------------------------------------
\begin{document}

% non-numbered pages
\pagestyle{empty} 

%----------------------------------------------------------------------------------------
%	TITLE
%----------------------------------------------------------------------------------------

% \begin{tabularx}{\linewidth}{ @{}X X@{} }
% \huge{Your Name}\vspace{2pt} & \hfill \emoji{incoming-envelope} email@email.com \\
% \raisebox{-0.05\height}\faGithub\ username \ | \
% \raisebox{-0.00\height}\faLinkedin\ username \ | \ \raisebox{-0.05\height}\faGlobe \ mysite.com  & \hfill \emoji{calling} number
% \end{tabularx}

\begin{tabularx}{\linewidth}{@{} C @{}}
\Huge{Jacques Reddinger} \\[7.5pt]
\href{https://github.com/jareddi}{\raisebox{-0.05\height}\faGithub\ jareddi} \ $|$ \ 
\href{https://linkedin.com/in/jacques-a-reddinger}{\raisebox{-0.05\height}\faLinkedin\ jacques-a-reddinger} \ $|$ \ 
\href{https://cas.uoregon.edu/directory/natural-sciences/all/jacquesr}{\raisebox{-0.05\height}\faGlobe \ University of Oregon} \ $|$ \ 
\href{mailto:jacques@reddinger.org}{\raisebox{-0.05\height}\faEnvelope \ jacques@reddinger.org} \\
% \href{tel:4433707163}{\raisebox{-0.05\height}\faMobile \ (443)370-7163} \\
\end{tabularx}

%----------------------------------------------------------------------------------------
% EXPERIENCE SECTIONS
%----------------------------------------------------------------------------------------

%Interests/ Keywords/ Summary
\section{Summary}
PhD researcher specializing in Electron Microscopy and Condensed Matter Physics with expertise in magnetic-enabled electron microscopy and micromagnetic simulation paired with custom data analysis tools. Published in peer-reviewed journals, presented at national and international conferences, and recognized with awards for research and teaching excellence.

%Experience
\section{Work Experience}

\begin{jobshort}{Graduate Research Assistant}{Sept 2020 - present}
\textit{McMorran Lab, Dept. of Physics, University of Oregon, Eugene, OR} \newline
Conducted experimental and computational studies of magnetic textures in magnetic multilayer structures. Designed and carried out experiments using advanced electron microscopy techniques. Reinforced experimental results with computational techniques. Collaborated with labs at University of California San Diego’s Center for Memory and Recording Research and the National Center for Electron Microscopy at Lawrence Berkeley National Lab. Co-authored 2 peer review papers and presented at 5 international conferences. Performed regular physics outreach.
\end{jobshort}

\begin{jobshort}{Undergraduate Research Assistant}{May 2019 - Aug 2020}
\textit{Koeth Lab, Dept. of Physics, University of Maryland, College Park, MD} \newline
Designed and conducted an experiment studying decay via electron capture of Be-7 ions. Restored the Small Isochronous Ring (SIR), an ion storage ring, to a functioning condition. Designed and implemented various critical support systems including power control, vacuum, cooling, and structural support.
\end{jobshort}

\begin{jobshort}{Undergraduate Research Assistant}{May 2019 - Aug 2020}
\textit{Belloni Lab, Dept. of Physics, University of Maryland, College Park, MD} \newline
Assisted in the search for the W$\gamma$ decay of a charged Higgs boson by analyzing high-energy physics data using ROOT, Python, and C++.
\end{jobshort}
  
%Projects
\section{Projects}

\begin{tabularx}{\linewidth}{ @{}l r@{} }
\textbf{Domain Tracker} & \hfill \href{https://github.com/jareddi/domain-tracker}{Link to github} \\[3.75pt]
\multicolumn{2}{@{}X@{}}{Custom python package for processing the output of micromagnetic simulation software. This tool harvests perpendicular magnetic domains and their properties from standard \texttt{.ovf} files. Domains are chronologically linked into \texttt{DomainChron} objects which store morphological, topological, and energetic properties. The information harvested and structured by this tool greatly assists in the analysis of simulation results.}  \\
\end{tabularx}

%----------------------------------------------------------------------------------------
%	EDUCATION
%----------------------------------------------------------------------------------------
\section{Education}
\begin{tabularx}{\linewidth}{@{}l X@{}}	
2020 - present & PhD Physics at \textbf{University of Oregon} \hfill \normalsize (GPA: 3.96/4.0) \\

2016 - 2020 & BS Physics at \textbf{University of Maryland} \hfill (GPA: 3.84/4.0) \\ 
\end{tabularx}

%----------------------------------------------------------------------------------------
%	PUBLICATIONS
%----------------------------------------------------------------------------------------
\section{Publications}
\begin{refsection}[cv.bib]
\nocite{*}
\newrefcontext[sorting=ydnt] % set reverse-chronological order
\printbibliography[heading=none]
\end{refsection}

%----------------------------------------------------------------------------------------
%	Conference Presentations & Posters
%----------------------------------------------------------------------------------------
\section{Conference Presentations \& Posters}
70th Annual Conference on Magnetism and Magnetic Materials, Palm Beach, FL, 31 October, 2025: "Soliton Formation and Topological Transitions during Stripe Collapse in Fe/Gd Multilayers"

Microscopy and Microanalysis 2025, Salt Lake City, UT, 30 July, 2025: "Dual-Modal Depth Sensitive Real Space Magnetic Imaging of 3D Skyrmions in Multilayer Thin Films"

16th Joint Conference on Magnetism and Magnetic Materials and Intermag, New Orleans, LA, 15 January, 2025: "Formation Process of Antiskyrmions in Fe/Gd Multilayers"

68th Annual Conference on Magnetism and Magnetic Materials, Dallas, TX, 1 November, 2023: "3D Dipole Antiskyrmions in Fe/Gd Multilayers Consist of Bloch Lines without Singularities"

%OMQ Fall Research Symposium 2023, Eugene, OR, September 11, 2023: ""

Microscopy and Microanalysis 2022, Portland, OR, 2 August, 2022: "Micromagnetics Simulation as a Supplement to and Diagnostic for Lorentz Transmission Electron Microscopy"

%----------------------------------------------------------------------------------------
%	AWARDS
%----------------------------------------------------------------------------------------
\section{Awards}
\begin{tabularx}{\linewidth}{@{}l X@{}}	
2025 & Science Graduate Student Research Program Awardee, Department of Energy\\
2025 & Microscopy and Microanalysis Student Scholar Award, Microanalysis Society \\

2022 & Best Student Poster Award, Microscopy Society of America \\

2022 & Weiser Physics Teaching Assistant Award, University of Oregon Dept. of Physics \\

2016 - 2020 & President’s Scholarship, University of Maryland \\ 
\end{tabularx}

%----------------------------------------------------------------------------------------
%	SKILLS
%----------------------------------------------------------------------------------------
\section{Technical Proficiencies}
\begin{tabularx}{\linewidth}{@{}l X@{}}
Electron Microscopy &  \normalsize{SEM, SEMPA, TEM, LTEM, STEM}\\
Surface Analysis &  \normalsize{plasma cleaning, ion milling, Auger spectroscopy, XRD}\\
Nanofabrication &  \normalsize{FIB, sputter deposition}\\
Experimental Apparatus &  \normalsize{UHV, electron/ion optics and detectors, power and measurement electronics}\\
Programming &  \normalsize{Python, Javascript, image analysis, computer vision}\\
Simulation & \normalsize{MuMax3, custom Monte-Carlo simulations}
\end{tabularx}

\vfill
\center{\footnotesize Last updated: \today}

\end{document}
